\documentclass[11pt,a4paper]{article}

% ==========================================
% 1. Encoding & Language
% ==========================================
\usepackage[utf8]{inputenc}
\usepackage[T1]{fontenc}
\usepackage[french]{babel}

% ==========================================
% 2. Modern Typography & Layout
% ==========================================
\usepackage{microtype}              % Subtle character protrusion & optical margin alignment
\usepackage[margin=2.5cm]{geometry} % Balanced margins

% ==========================================
% 3. Core Math & Utility Packages
% ==========================================
\usepackage{amsmath,amsthm}
\usepackage{libertine}
\usepackage[libertine]{newtxmath} % Automatically imports AMS math symbols
\usepackage{enumitem}
\usepackage{tikz}
\usepackage[most]{tcolorbox}
\usepackage{xcolor}
\usepackage{caption}


% Configure compact, professional lists
\setlist[enumerate]{noitemsep, topsep=3pt, parsep=0pt, partopsep=0pt}
\setlist[itemize]{noitemsep, topsep=3pt, parsep=0pt, partopsep=0pt}

% ==========================================
% 4. Custom Colors & Visual Enhancements
% ==========================================
\definecolor{primary}{RGB}{37, 99, 235}      % Modern slate blue
\definecolor{secondary}{RGB}{100, 116, 139}  % Subtle border grey
\definecolor{boxbg}{RGB}{248, 250, 252}      % Soft neutral background

% ==========================================
% 5. Modern Card Environments (tcolorbox)
% ==========================================
% Accent command for terms
\newcommand{\defemph}[1]{\textbf{\color{primary}#1}}
\newcommand{\ol}[1]{\overline{#1}}

% Clean Left-Border Exercise Box
\newtcolorbox{td-exo}[1][]{
    enhanced,
    breakable,
    colback=boxbg,
    colframe=primary,
    boxrule=0pt,
    leftrule=3.5pt,
    arc=0mm,
    top=8pt, bottom=8pt, left=12pt, right=12pt,
    title={\bfseries\color{primary}Exercice : #1},
    detach title,
    before upper={\tcbtitle\par\vspace{4pt}},
    after title={\par},
    fonttitle=\bfseries
}

% Solution Container
\newtcolorbox{td-sol}[1][]{
    enhanced,
    breakable,
    colback=white,
    colframe=secondary,
    boxrule=0.5pt,
    arc=1mm,
    top=8pt, bottom=8pt, left=12pt, right=12pt,
    title={\bfseries\color{secondary}Solution #1},
    detach title,
    before upper={\tcbtitle\par\vspace{4pt}},
    fonttitle=\bfseries
}

% Theorem / Definition Styles (amsthm with break hook)
\newtheoremstyle{break}%
  {6pt}{6pt}%
  {\normalfont}{\parindent}%
  {\bfseries\color{primary}}{.}%
  {\newline}{}%

\theoremstyle{break}
\newtheorem{definition}{Définition}
\newtheorem{theorem}{Théorème}

% ==========================================
% 6. Document Metadata
% ==========================================
\title{\textbf{HAI9151 : Graphes, Complexité et Algorithmes}\\[0.2em]\large Complément de TD — Mineurs de graphes}
\author{\textbf{Ivan Lejeune}}
\date{24 septembre 2026}

% ==========================================
% DOCUMENT BODY
% ==========================================
\begin{document}

\maketitle
\tableofcontents
\vspace{1.5em}

\section{Mineurs de graphes}\label{sec:mineurs-de-graphes}

% --- EXERCISE 1 ---
\begin{td-exo}[Ordres graphiques]
    \begin{enumerate}
        \item Montrer que la relation de sous-graphe forme un ordre sur les graphes.
        \item On définit la relation suivante sur les graphes :
        \begin{equation*}
            G \ll G' \quad\text{si}\quad n(G) \le n(G') \quad\text{et}\quad m(G) \le m(G').
        \end{equation*}
        La relation \(\ll\) forme-t-elle un ordre sur les graphes ?
    \end{enumerate}
\end{td-exo}

\begin{td-sol}
    \begin{enumerate}
        \item On rappelle la définition d'un ordre :
        
        \begin{definition}
            Soit \(E\) un ensemble et \(\mathcal{R}\) une relation entre deux éléments de \(E\).
            On dit que \(\mathcal{R}\) est une \defemph{relation d'ordre} sur \(E\) si elle vérifie les propriétés suivantes :
            \begin{enumerate}[label=\roman*)]
                \item Réflexivité : pour tout \(x \in E\), on a \(x\mathcal{R}x\).
                \item Antisymétrie : pour tous \(x,y\in E\), si \(x\mathcal{R}y\) et \(y\mathcal{R}x\) alors \(x = y\).
                \item Transitivité : pour tous \(x,y,z\in E\), si \(x\mathcal{R}y\) et \(y\mathcal{R}z\) alors \(x\mathcal{R}z\).
            \end{enumerate}
        \end{definition}

        Vérifions alors ces propriétés sur la relation de sous-graphe :
        \begin{enumerate}[label=(\alph*)]
            \item \textbf{Réflexivité :} pour tout graphe \(G\), on a bien \(V(G) \subseteq V(G)\) et \(E(G) \subseteq E(G)\), donc \(G\) est bien sous-graphe de \(G\).

            \item \textbf{Antisymétrie :} soient \(G\) et \(H\) deux graphes tels que \(G\) est sous-graphe de \(H\) et \(H\) est sous-graphe de \(G\).
            \begin{equation*}
                \begin{aligned}
                    V(G) \subseteq V(H) \quad\text{et}\quad V(H) \subseteq V(G) &\implies V(G) = V(H) \\
                    E(G) \subseteq E(H) \quad\text{et}\quad E(H) \subseteq E(G) &\implies E(G) = E(H)
                \end{aligned}
            \end{equation*}
            Ainsi, comme \(V(G) = V(H)\) et \(E(G) = E(H)\), on a bien \(G = H\).

            \item \textbf{Transitivité :} soient \(G, H\) et \(I\) trois graphes tels que \(G\) est sous-graphe de \(H\) et \(H\) est sous-graphe de \(I\).
            \begin{equation*}
                \begin{aligned}
                    V(G) \subseteq V(H) \quad\text{et}\quad V(H) \subseteq V(I) &\implies V(G) \subseteq V(I) \\
                    E(G) \subseteq E(H) \quad\text{et}\quad E(H) \subseteq E(I) &\implies E(G) \subseteq E(I)
                \end{aligned}
            \end{equation*}
            Ainsi, comme \(V(G) \subseteq V(I)\) et \(E(G) \subseteq E(I)\), on a bien \(G\) sous-graphe de \(I\).
        \end{enumerate}

        \item On pose \(G = ([a, b, c, d], [(a, b), (c, d)])\) et \(H = ([a, b, c, d], [(a, b), (b, c)])\) comme dans la figure ci-dessous :
        
        \vspace{0.5em}
        \begin{center}
            \begin{center}
\tikzset{
    main node/.style={
        circle, 
        draw, 
        fill=blue!20, 
        inner sep=1pt, 
        font=\scriptsize, 
        minimum size=6mm, 
        text=black
    }
}

% Graph (a)
\begin{minipage}[t]{0.45\linewidth}
    \centering
    \begin{tikzpicture}[scale=1.5]
        \node[main node] (1) at (0,0) {\(a\)};
        \node[main node] (2) at (1,0) {\(b\)};
        \node[main node] (3) at (0,-1) {\(c\)};
        \node[main node] (4) at (1,-1) {\(d\)};
        \draw (1) -- (2);
        \draw (3) -- (4);
    \end{tikzpicture}
    
    \vspace{0.4em}
    {\small (a) Graphe \(G\)}
\end{minipage}%
\hfill
% Graph (b)
\begin{minipage}[t]{0.45\linewidth}
    \centering
    \begin{tikzpicture}[scale=1.5]
        \node[main node] (1) at (0,0) {\(a\)};
        \node[main node] (2) at (1,0) {\(b\)};
        \node[main node] (3) at (0,-1) {\(c\)};
        \node[main node] (4) at (1,-1) {\(d\)};
        \draw (1) -- (2);
        \draw (2) -- (3);
    \end{tikzpicture}
    
    \vspace{0.4em}
    {\small (b) Graphe \(H\)}
\end{minipage}

\vspace{0.8em}
\captionof{figure}{Illustration des graphes \(G\) et \(H\)}\label{fig:td1_ex1_fi1}
\end{center}
            \captionof{figure}{Illustration des graphes \(G\) et \(H\)}
        \end{center}
        \vspace{0.5em}

        On a \(n(G) = n(H)\) et \(m(G) = m(H)\) donc \(G \ll H\) et \(H \ll G\).
        Pourtant, \(G \neq H\), donc la propriété d'antisymétrie n'est pas respectée et \(\ll\) ne définit pas un ordre sur les graphes.
    \end{enumerate}
\end{td-sol}

% --- EXERCISE 2 ---
\begin{td-exo}[Sous-graphe induit dense]
    Montrer que la borne donnée par le théorème du sous-graphe induit dense est optimale.
\end{td-exo}

\begin{td-sol}
    On rappelle le théorème :
    \begin{theorem}[Sous-graphe induit dense]
        Si le degré moyen d'un graphe \(G\) est \(\geq 2k\), alors \(G\) contient un sous-graphe induit \(F\) de degré minimum \(\geq k + 1\).
    \end{theorem}

    Soit \(G\) un cycle de longueur \(n\). Alors, on a \(\ol{d}(G) = 2k = 2 \times 1\).
    On peut bien trouver un sous-graphe induit \(F\) de degré minimum \(\geq k + 1 = 2\) (par exemple, le cycle lui-même) mais il n'existe pas de sous-graphe induit \(F\) de degré minimum \(\geq k + 2 = 3\), car tous les sommets du cycle ont un degré de \(2\).

    La borne est donc optimale.

    On peut aussi construire un exemple plus général. On construit \(G\) de la manière suivante :
    \begin{itemize}
        \item On prend un graphe complet \(K_{k + 1}\) de degré moyen \(k\).
        \item On ajoute au graphe un stable de taille \(M\) et on crée le biparti complet entre le stable et le graphe complet.
    \end{itemize}
    On obtient alors un graphe \(G\) de degré total
    \begin{equation*}
        (k+1)(k+M) + M(k+1) = (k+1)(k+2M)
    \end{equation*}
    et de degré moyen
    \begin{equation*}
        \frac{(k+1)(k+2M)}{k+1+M}.
    \end{equation*}
    En développant, on trouve que pour \(M \geq \frac{k^2 + k}{2}\), on a bien \(\ol{d}(G) \geq 2k\).
    Comme le graphe complet \(K_{k + 1}\) est un sous-graphe induit de \(G\), on a bien un sous-graphe induit de degré minimum \(\geq k + 1\), mais il n'existe pas de sous-graphe induit de degré minimum \(\geq k + 2\), car tous les sommets du stable ont un degré de \(k + 1\).
\end{td-sol}

% --- EXERCISE 3 ---
\begin{td-exo}[Gros sous-graphe couvrant biparti]
    Montrer que tout graphe \(G\) contient un sous-graphe biparti couvrant \(H\) avec \(m(H) \geq m(G)/2\) (on pourra ordonner les sommets de \(G\) en \(v_1,\ldots,v_n\) et construire la bipartition \(X, Y\) de \(H\) en plaçant \(v_i\) dans \(X\) ou \(Y\) en fonction des \(v_j\) pour \(j < i\)).
\end{td-exo}

\begin{td-sol}
    Soit \(H\) le sous-graphe induit de \(G\) de degré moyen maximum et de nombre de sommets minimum pour cela.

    Si \(\delta(H) \geq k + 1\), la condition est directement vérifiée. Sinon, \(H\) contient un sommet \(x\) tel que \(d_H(x) \leq k\).

    En particulier, on a \(\ol{d}(H) \geq \ol{d}(G) \geq 2k\). On déduit :
    \begin{equation*}
        \begin{aligned}
            2m(H \setminus x) &= 2m(H) - 2d_H(x) \geq 2(m(H) - k) \\
            &\geq \frac{2}{n(H)} \cdot \left(n(H) m(H) - n(H) k\right) \\
            &\geq \frac{2}{n(H)} \cdot \left(n(H) m(H) - m(H)\right) \\
            &\geq 2 \cdot \frac{m(H)}{n(H)} \cdot (n(H) - 1) \\
            &= \ol{d}(H) \cdot n(H \setminus x).
        \end{aligned}
    \end{equation*}
    Et donc :
    \begin{equation*}
        \ol{d}(H \setminus x) = 2 \frac{m(H \setminus x)}{n(H \setminus x)} \geq \ol{d}(H).
    \end{equation*}

    Si \(\ol{d}(H \setminus x) > \ol{d}(H)\), cela contredit le choix de \(H\).
    Si \(\ol{d}(H \setminus x) = \ol{d}(H)\), cela contredit également le choix de \(H\) car \(H \setminus x\) possède un sommet de moins que \(H\).
\end{td-sol}

\end{document}